\documentclass[11pt,a4paper]{article}
\usepackage[a4paper,margin=27mm,headheight=14pt]{geometry}
\usepackage{fontspec}
\setmainfont{Latin Modern Roman}
\setsansfont{Latin Modern Sans}
\setmonofont{Latin Modern Mono}
\usepackage{amsmath,amssymb,amsthm,mathtools,mathrsfs,enumitem,booktabs,array,longtable,fancyhdr}
\usepackage{xcolor}
\definecolor{MathInternalLink}{HTML}{1F4E79}
\definecolor{MathSectionLink}{HTML}{00656D}
\definecolor{MathEquationLink}{HTML}{643F7A}
\definecolor{MathCitationLink}{HTML}{2F6B3C}
\definecolor{MathExternalLink}{HTML}{843C39}
\usepackage{hyperref}
\hypersetup{colorlinks=true,linkcolor=MathInternalLink,citecolor=MathCitationLink,
urlcolor=MathExternalLink,linktoc=all,bookmarksnumbered=true,
pdftitle={Arithmetic Conjectures for Beilinson--Bloch Heights on Products of Curves},
pdfauthor={}}
\newtheorem{theorem}{Theorem}[section]
\newtheorem{proposition}[theorem]{Proposition}
\newtheorem{lemma}[theorem]{Lemma}
\newtheorem{corollary}[theorem]{Corollary}
\theoremstyle{definition}
\newtheorem{definition}[theorem]{Definition}
\newtheorem{conjecture}[theorem]{Conjecture}
\newtheorem{example}[theorem]{Example}
\newtheorem{question}[theorem]{Question}
\usepackage[nameinlink,noabbrev]{cleveref}
\crefname{conjecture}{conjecture}{conjectures}
\Crefname{conjecture}{Conjecture}{Conjectures}
\crefname{question}{question}{questions}
\Crefname{question}{Question}{Questions}
\DeclareRobustCommand{\MathSecref}[1]{\begingroup\hypersetup{linkcolor=MathSectionLink}\Cref{#1}\endgroup}
\DeclareRobustCommand{\MathEqref}[1]{\begingroup\hypersetup{linkcolor=MathEquationLink}\eqref{#1}\endgroup}
\newcommand{\Q}{\mathbb Q}\newcommand{\R}{\mathbb R}\newcommand{\C}{\mathbb C}
\newcommand{\Z}{\mathbb Z}\newcommand{\A}{\mathbb A}\newcommand{\Fbar}{\overline F}
\newcommand{\CH}{\operatorname{CH}}\newcommand{\Pic}{\operatorname{Pic}}
\newcommand{\Hom}{\operatorname{Hom}}\newcommand{\End}{\operatorname{End}}
\newcommand{\Spec}{\operatorname{Spec}}\newcommand{\Gal}{\operatorname{Gal}}
\newcommand{\rad}{\operatorname{rad}}\newcommand{\rk}{\operatorname{rank}}
\newcommand{\im}{\operatorname{im}}\newcommand{\Griff}{\operatorname{Griff}}
\newcommand{\AJ}{\operatorname{AJ}}\newcommand{\Tr}{\operatorname{Tr}}
\newcommand{\cl}{\operatorname{cl}}\newcommand{\ch}{\operatorname{ch}}
\newcommand{\Td}{\operatorname{Td}}\newcommand{\Supp}{\operatorname{Supp}}
\newcommand{\homol}{\mathrm{hom}}\newcommand{\alg}{\mathrm{alg}}
\newcommand{\et}{\textnormal{ét}}\newcommand{\cont}{\mathrm{cont}}
\newcommand{\cP}{\mathcal P}\newcommand{\cO}{\mathcal O}
\newcommand{\cD}{\mathscr D}\newcommand{\cV}{\mathcal V}
\newcommand{\tGamma}{{}^t\Gamma}
\numberwithin{equation}{section}
\setlist{leftmargin=2em,itemsep=3pt,topsep=5pt}
\setlength{\emergencystretch}{2em}
\setlength{\parskip}{3pt}
\pagestyle{fancy}\fancyhf{}
\fancyhead[L]{\small Arithmetic conjectures for cycle heights}
\fancyfoot[C]{\thepage}
\title{Arithmetic Conjectures for Beilinson--Bloch Heights\\
on Products of Curves}
\author{}
\date{17 November 2026}
\begin{document}
\maketitle
\begin{abstract}
We formulate an arithmetic programme for height pairings on the
homologically trivial Chow groups of products of curves. Exact
Chow projectors reduce arithmetic hard Lefschetz to off-central
vanishing and reduce primitive integral Northcott to two central
assertions: injectivity of the algebraic Abel--Jacobi map and
Northcott on a residual Griffiths quotient. The latter assertion
contains real positivity but also requires integral and
varying-field information. We separate the proposed arithmetic
inputs into scalar detection by continuous étale cohomology,
a common rational lattice, and ramification control for the
minimal fields of cycle classes. The stated implications specify
where each input enters.
\end{abstract}
{\hypersetup{linkcolor=MathSectionLink}\tableofcontents}

\section{Introduction}\label{sec:introduction}

The expected positivity of Beilinson--Bloch heights is an arithmetic
Hodge--Riemann principle. A Northcott property is stronger: it asks
for finiteness of integral cycle classes of bounded degree and
height. Even a positive quadratic form does not give this
finiteness unless the integral classes are arithmetically discrete.
Allowing the field of definition to vary adds a further descent
problem.

For products of curves, tensor Chow projectors separate these
questions. The all-\(h^1\) part of a subproduct carries a central
height precisely when its dimension is \(2c-1\) in codimension
\(c\). The other parts are responsible for arithmetic Lefschetz.
Within the central part, curve correspondences separate a
Néron--Tate summand and its residual Griffiths quotient. This
paper formulates the resulting conjectures and organizes the
arithmetic inputs to their study. The complete reduction proofs
are given in the companion manuscript \cite{ArithmeticReductions}.

These conjectures are expressed within the Bloch--Beilinson
and arithmetic standard-conjecture framework. The arithmetic
Lefschetz and primitive positivity statements are those recorded
in \cite[Conjecture~2.3.1]{ZhangStandard}. The purpose of the
formulation here is to specify the residual objects, their
integral images and the exact implications connecting the
different arithmetic questions.

\subsection{The fixed product and height data}\label{sec:height-data}

Let \(F\) be a number field, and let
\[
 X=\prod_{i=1}^n C_i,\qquad
 o_i\in C_i(F),\qquad H_i=p_i^*[o_i],\qquad
 \ell=\sum_{i=1}^n a_iH_i,\quad a_i\in\Z_{>0},
\]
where the curves are smooth, projective and geometrically
connected. For a finite extension \(E/F\), write
\[
 V^r(X_E)=\CH^r(X_E)_{\homol,\Q}.
\]

% Source: papers/arithmetic-sections/conventions.tex
Fix an algebraic closure \(\Fbar\) of \(F\) and an embedding into
\(\C\). Homological triviality is defined by geometric rational
Betti cohomology. All Chow groups below have rational coefficients
unless otherwise indicated. We regard \(V^r(X_E)\) as a subspace
of \(V^r(X_{\Fbar})\): restriction is injective because a rational
equivalence descends to a finite extension and trace composed with
restriction multiplies by its degree.

\begin{definition}\label[definition]{def:heights}
A curve-normalized height system consists of pairings on all smooth
projective varieties used below, including the subproducts and the
source curves of correspondences. For such a variety \(X\) of
dimension \(n\), these are
\[
 h_X:V^r(X_E)\times V^{n+1-r}(X_E)\longrightarrow\R
\]
over every finite extension \(E/F\), satisfying:
\begin{enumerate}[label=(H\arabic*)]
\item bilinearity over \(\Q\) and symmetry under exchange of the
complementary degrees;
\item invariance under extension in the absolute normalization,
and covariance under isomorphisms, including Galois conjugation;
\item correspondence adjunction
\[
 h_X(\Gamma_*z,w)=h_Y(z,{}^t\Gamma_*w);
\]
\item on a pointed curve \(C\), the form \(-h_C\) on
\(\Pic^0(C)_\Q\) is the positive Néron--Tate form for the canonical
principal polarization, in the corresponding Poincaré normalization.
\end{enumerate}
\end{definition}

These hypotheses will be understood in every statement involving
\(h\). A real positive form means positive definiteness after
tensoring the rational group over a fixed number field with \(\R\).
For codimension \(r\) on \(X\), the symbol \(\AJ_\Q\)
denotes the complex Abel--Jacobi map with target
\(J^r(X_\C)(\C)\otimes_\Z\Q\).

Let \(\Lambda_X^r(E)\) be the image in \(V^r(X_{\Fbar})\) of the
integral-coefficient cycles over \(E\) with zero geometric rational
cycle class. Thus torsion and rational equivalence are removed before
counting. For a fixed rational projector \(P\), we use its projected
integral image \(P_*\Lambda_X^r(E)\); for an integral projector this
equals \(\Lambda_X^r(E)\cap P_*V^r(X_E)\).

\begin{definition}\label[definition]{def:northcott}
For compatible integral images \(\Lambda(E)\) in one geometric
rational group and a quadratic form \(q\), bounded-degree Northcott
means that, for every \(B\in\R\) and integer \(D\ge1\),
\[
 \{\xi\in\bigcup_{[E:F]\le D}\Lambda(E):q(\xi)\le B\}
\]
is finite. Fixed-field Northcott means the same assertion with
\(E\) fixed. Bounded denominators mean the union of
\(d^{-1}\Lambda(E)\) for \(1\le d\le d_0\), with \(d_0\) fixed.
\end{definition}

Northcott for integral images implies the bounded-denominator version
by scaling. A fixed initial extension \(F_0/F\) is harmless for
bounded-degree assertions: after replacing \(E\) by \(EF_0\),
\([EF_0:F_0]\le[E:F]\), the integral image remains integral, and
the absolute height is unchanged.





\subsection{Central cores and the residual pairing}
\label{sec:central-data}

For a subset \(S\subseteq\{1,\ldots,n\}\), put
\[
 Y_S=\prod_{i\in S}C_i,\qquad
 \Pi_S=\bigotimes_{i\in S}\pi_{1,i},
 \qquad U_S^c(E)=\Pi_{S*}V^c(Y_{S,E}).
\]
Here \(\pi_{1,i}=\Delta_{C_i}-o_i\times C_i-C_i\times o_i\).
The empty subproduct has zero homologically trivial Chow groups.
The core is central if \(|S|=2c-1\).

Fix such an odd subproduct. Let
\[
 A_{S,E}=U_S^c(E)\cap\CH^c(Y_{S,E})_{\alg,\Q},
 \qquad G_{S,E}=U_S^c(E)/A_{S,E}.
\]
The maximal curve-visible construction in
\cite[Section 3]{ArithmeticReductions} gives an actual
self-transpose Chow projector \(e_S\le\Pi_S\), with
\(R_S=\Pi_S-e_S\). Choose one finite extension \(F_0/F\)
over which these projectors are defined for all odd subproducts.
For every finite \(E/F_0\),
\[
 U_S^c(E)=e_SU_S^c(E)\oplus R_SU_S^c(E),\qquad
 e_SU_S^c(E)\simeq B_S(E)_\Q
\]
for a fixed abelian variety \(B_S/F_0\).
The decomposition is height-orthogonal, and the form
\(q_{e,S}(u)=(-1)^ch_{Y_S}(u,u)\) on the first summand is
a positive Néron--Tate form. Moreover,
\[
 R_SA_{S,E}=\ker(\AJ_\Q|_{A_{S,E}})
 =A_{S,E}\cap\rad(h_{Y_S}|_{U_S^c(E)}).
\]
For \(\xi=[z]\), \(\eta=[w]\) in \(G_{S,E}\), define
\[
 h_{G,S}(\xi,\eta)=h_{Y_S}(R_Sz,R_Sw),\qquad
 q_{G,S}(\xi)=(-1)^ch_{G,S}(\xi,\xi).
\]
The radical equality makes this pairing well defined, and the
projector construction proves its independence of the maximal
curve-visible lift. These facts hold for every height system in
\Cref{def:heights}.

The integral structures are
\[
 \Lambda_{U,S}(E)=\Pi_{S*}\Lambda_{Y_S}^c(E),\qquad
 \Lambda_{G,S}(E)=\im\!\left(\Lambda_{U,S}(E)\to G_{S,E}\right).
\]
The projector \(\Pi_S\) is integral. The residual projector
\(R_S\) is rational, with one fixed denominator; that denominator
is independent of the cycle and of \(E\).
Since algebraic triviality is geometric, all these rational
quotients embed in their counterparts over \(\Fbar\).

\subsection{The three central assertions}\label{sec:three-assertions}

\begin{conjecture}[Off-central vanishing]
\label[conjecture]{conj:offcentral}
For every subproduct \(Y_S\), every finite \(E/F_0\), and
every \(c\) with \(|S|\ne2c-1\),
\[
 U_S^c(E)=0.
\]
\end{conjecture}

\begin{conjecture}[Central algebraic Abel--Jacobi injectivity]
\label[conjecture]{conj:AJ}
For every \(|S|=2c-1\) and finite \(E/F_0\), the map
\[
 \AJ_\Q:A_{S,E}\longrightarrow
 J^c(Y_{S,\C})(\C)\otimes_\Z\Q
\]
is injective.
\end{conjecture}

\begin{conjecture}[Residual integral Northcott]
\label[conjecture]{conj:residual-northcott}
For every \(|S|=2c-1\), every \(B\in\R\), and every integer
\(D\ge1\), the set of geometric classes
\[
 \left\{\xi\in\bigcup_{[E:F_0]\le D}\Lambda_{G,S}(E):
       q_{G,S}(\xi)\le B\right\}
\]
is finite.
\end{conjecture}

These assertions address, respectively, the degrees in which
the height pairing should be nondegenerate, its radical on
the algebraically trivial part, and the arithmetic size of
the remaining integral quotient.

\section{From the central assertions to the whole product}
\label{sec:product-programme}

\subsection{Arithmetic Lefschetz and the primitive form}
\label{sec:lefschetz-conjecture}

For \(1\le r\le\lfloor(n+1)/2\rfloor\), put
\[
\begin{aligned}
 P_\ell^r(X_E)&=\ker\!\left(
   \ell^{n+2-2r}:V^r(X_E)\longrightarrow V^{n+2-r}(X_E)\right),\\
 q_{\ell,r}(z)&=(-1)^rh_X(z,\ell^{n+1-2r}z).
\end{aligned}
\]
The primitive integral image is
\(\Lambda_X^r(E)\cap P_\ell^r(X_E)\).

\begin{theorem}[Equivalence of the arithmetic programmes]
\label[theorem]{thm:program-equivalence}
For the fixed pointed product, positive integers \(a_i\), and
a height system satisfying \Cref{def:heights}, the following
conditions are equivalent.
\begin{enumerate}[label=(\roman*)]
\item \Cref{conj:offcentral,conj:AJ,conj:residual-northcott}
hold for all subproducts of \(X\).
\item For every finite \(E/F_0\) and every
\(1\le r\le\lfloor(n+1)/2\rfloor\),
\[
 \ell^{n+1-2r}:V^r(X_E)\xrightarrow{\sim}V^{n+1-r}(X_E),
\]
and \(q_{\ell,r}\) has bounded-degree Northcott on the
primitive integral images.
\end{enumerate}
Either condition implies real positive definiteness of all
primitive forms and nondegeneracy, also after real scalar
extension, of the complementary height pairings over every
fixed number field.
\end{theorem}

The full proof is the principal reduction theorem of
\cite{ArithmeticReductions}, completed in its Section 7.
Its two mechanisms are the exact central-block identity and
the central radical decomposition. We recall them because
they determine which conjecture is needed at each step.

The Chow block indexed by \(S,c\) is
\[
 U_S^c(E)\otimes_\Q \Q[x_i:i\notin S]/(x_i^2),
\]
with \(\ell\) acting as multiplication by
\(\sum_{i\notin S}x_i\). Writing
\(I=\{1,\ldots,n\}\setminus S\), the identification sends
\(u\otimes x_T\) to \(a_Tj_Tu\), where
\(a_T=\prod_{i\in T}a_i\) and \(j_T\) inserts \([o_i]\)
on \(T\) and \([C_i]\) on \(I\setminus T\). The lowering
operator is
\[
 F x_T=\sum_{i\in T}x_{T\setminus\{i\}}.
\]
The arithmetic hard Lefschetz maps
are isomorphisms exactly when the noncentral cores vanish.
This equivalence is independent of the choice of the positive
weights \(a_i\), with the points \(o_i\) fixed.

In a central block let \(m=n-|S|\), \(t=r-c\),
\(I=\{1,\ldots,n\}\setminus S\), and
\(z_S=\sum_{|T|=t}u_Tx_T\). For \(0\le t\le m/2\),
primitivity means \(Fz_S=0\); the signed height is then
\begin{equation}\label{eq:programme-squares}
 q_{\ell,r}(z_S)
  =(m-2t)!\prod_{i\in I}a_i
       \sum_{|T|=t}(-1)^c h_{Y_S}(u_T,u_T).
\end{equation}
This central-block identity uses only adjunction and the
square-zero calculation. Off-central vanishing enters when
all of \(V^r(X_E)\) is reduced to these central blocks.
If \(z\) is integral, the single integer
\(N_a=\prod_i a_i\) clears the denominators of every \(u_T\).

Finally, central integral Northcott is equivalent to
\[
 \ker(\AJ_\Q|_{A_{S,E}})=0\quad\hbox{for every finite }E/F_0,
 \qquad\text{and residual integral Northcott on }\Lambda_{G,S}.
\]
Indeed
\[
 (-1)^ch_{Y_S}(z,z)
  =q_{e,S}(e_Sz)+q_{G,S}([z]).
\]
The visible summand has classical bounded-degree Northcott
through fixed correspondences with bounded denominators.
A nonzero algebraic radical class would have infinitely
many integral multiples of height zero. After this radical
is removed, the pair \((e_Sz,[z])\) detects \(z\).

\subsection{Residual positivity and its precise role}
\label{sec:residual-positivity}

\begin{conjecture}[Real residual positivity]
\label[conjecture]{conj:residual-positivity}
For every central subproduct and every finite \(E/F_0\),
\(q_{G,S}\) extends to a positive definite quadratic form on
\(G_{S,E}\otimes_\Q\R\).
\end{conjecture}

This is a weaker consequence of
\Cref{conj:residual-northcott}. More generally, if a subgroup
\(\Lambda\) spans a rational vector space \(V\), finiteness
of all height balls in \(\Lambda\) forces the real extension
of \(q\) to be positive definite. A negative direction gives
a negative rational vector and its integer multiples.
A real null direction, in a finite-dimensional rational
subspace, gives infinitely many integral vectors of bounded
height by rounding its scalar multiples. This is the proof
of the real-positivity lemma in
\cite[Section 5.1]{ArithmeticReductions}.

The converse is an integral problem. For example,
\[
 \Lambda=\sum_{p\ {\rm prime}}\Z\,p^{-1}\subset\Q,
 \qquad q(x)=x^2
\]
has a positive real form and a finite free local completion
\(\Lambda\otimes\Z_p=p^{-1}\Z_p\) at every prime. Nevertheless
\(q(1/p)\to0\). The extra input needed for finiteness is one
common rational lattice, rather than a separate local bound
at each prime.

With \Cref{conj:AJ,conj:residual-positivity}, the central form
is real positive definite by its orthogonal splitting.
Adding \Cref{conj:offcentral} propagates this positivity and
nondegeneracy to the product through
\eqref{eq:programme-squares}. The associated primitive
bilinear form then satisfies
\[
 |b_{\ell,r}(z,w)|^2\le q_{\ell,r}(z)q_{\ell,r}(w).
\]
The finiteness assertion in \Cref{conj:residual-northcott}
requires the additional arithmetic inputs below.

\section{Arithmetic detection and the integral lattice}
\label{sec:detection-conjectures}

Fix one central subproduct \(Y\) of dimension \(k=2c-1\),
and abbreviate \(G_{S,E},R_S,q_{G,S}\) to \(G_E,R,q_G\).

\subsection{Scalar detection by a residual regulator}
\label{sec:scalar-detection}

For every prime \(p\), let
\[
 V_p=R_*H_{\et}^{2c-1}(Y_{\Fbar},\Q_p(c)).
\]
The actual projector \(R\) has geometric cohomology in one
degree. Consequently, its part of Hochschild--Serre has one
row and gives a canonical identification
\[
 R_*H^{2c}_{\cont}(Y_E,\Q_p(c))\simeq H^1(E,V_p).
\]
The continuous cycle class of \(Rz\) annihilates \(A_E\).
It therefore defines a residual arithmetic Abel--Jacobi map
\[
 a_{p,E}:G_E\longrightarrow H^1(E,V_p),
\]
compatible with restriction and trace and independent of
the curve-visible lift under the common cohomological
identification. The construction is proved in
\cite[Section 4.1]{ArithmeticReductions}, in the
continuous-cohomology framework of \cite{Raskind,Kahn}.

The representation \(V_p\) has weight \(-1\). For one fixed
set \(S_0\) of places of \(F_0\), containing the bad-reduction
places, inflation identifies
\[
 H^1(G_{E,S_E},V_p)\simeq H^1(E,V_p),
 \qquad S_E=\{v:v\mid S_0\text{ or }v\mid p\},
\]
with archimedean places included. This whole target is
finite-dimensional over \(\Q_p\).

\begin{conjecture}[Residual scalar detection]
\label[conjecture]{conj:detection}
For every finite \(E/F_0\) and every prime \(p\), the
scalar-extended map
\[
 G_E\otimes_\Q\Q_p\longrightarrow H^1(E,V_p)
\]
is injective.
\end{conjecture}

The word \emph{scalar} specifies the strength needed here.
Injectivity of the \(\Q\)-linear map \(a_{p,E}\) alone would
not imply finite rational rank: a finite-dimensional
\(\Q_p\)-space is infinite-dimensional over \(\Q\).
In contrast, scalar-extended injectivity at one prime gives
\[
 \dim_\Q G_E
 \le \dim_{\Q_p}H^1(E,V_p)<\infty.
\]

For every fixed \(E,p\), the integral cycle map and one
denominator clearing \(R\) place
\(a_{p,E}(\Lambda_G(E))\) in a finitely generated
\(\Z_p\)-module. Under scalar-extended injectivity,
\[
 \Lambda_G(E)\otimes_\Z\Z_p
 \quad\text{is finitely generated over }\Z_p,
 \qquad \Lambda_G(E)/p\Lambda_G(E)\quad\text{is finite}.
\]
The proof uses the integral image theorem of
\cite[Theorem 2.4]{Kahn}; it is given in
\cite[Section 4.3]{ArithmeticReductions}.

\subsection{A common lattice across coefficient primes}
\label{sec:lattice-conjecture}

\begin{conjecture}[Residual integral discreteness]
\label[conjecture]{conj:lattice}
For every finite \(E/F_0\), the space \(G_E\) is
finite-dimensional, and there are a rational basis
\(v_1,\ldots,v_s\) and an integer \(N_E\ge1\) such that
\[
 \Lambda_G(E)\subseteq N_E^{-1}\sum_{i=1}^s\Z v_i.
\]
\end{conjecture}

For a finite-rank torsion-free group and a real positive
definite quadratic form, the following conditions are
equivalent: finite height balls, a positive height gap,
finite generation, a common denominator bound in a rational
basis, and relative compactness of the diagonal image in
the finite adeles. The proof in
\cite[Section 5.2]{ArithmeticReductions} is an elementary
lattice argument. It follows that
\Cref{conj:lattice,conj:residual-positivity} imply
fixed-field Northcott for every \(E\).
Conversely, fixed-field Northcott together with finite
rational rank implies \Cref{conj:lattice} and real
positivity. Finite rank is an explicit premise in this
converse.

A route to \Cref{conj:lattice} is an almost-everywhere
saturation comparison. Fix \(E\), put \(V=G_E\), and suppose
that \(V\) is finite-dimensional. Seek a full rational
lattice \(\Lambda_0\subset V\) and full local lattices
\(M_p\subset H^1(E,V_p)\) such that the scalar-extended
regulators \(D_p\) are injective and
\begin{equation}\label{eq:common-lattice-target}
 D_p(\Lambda_G(E))\subseteq M_p\quad\text{for every }p,
 \qquad
 D_p^{-1}(M_p)=\Lambda_0\otimes\Z_p
       \quad\text{for almost every }p.
\end{equation}
Then the inverse lattices at the finitely many exceptional
primes can be absorbed into one integer \(N_E\).
Thus \eqref{eq:common-lattice-target} proves
\Cref{conj:lattice}. The fixed-prime local bounds above
supply the first type of control; the equality with one
common rational lattice is the comparison still required.

Integral detection is formulated on the stated image in
a rational quotient. The unrestricted integral continuous
cycle class map has torsion counterexamples over number
fields in codimension three
\cite[Theorems 1.3--1.4]{ScaviaSuzuki}.
The torsion kernels in those statements and the torsion-free
group \(\Lambda_G(E)\) are different arithmetic objects.

\subsection{The height-gap route}\label{sec:height-gap-route}

There is a second route that can start from one coefficient
prime. Suppose \(\sqrt{q_G}\) is the norm of a real positive
quadratic form, and, for some integer \(m\ge2\),
\[
 s=|\Lambda_G(E)/m\Lambda_G(E)|<\infty,\qquad
 q_G(\lambda)\ge\varepsilon_E>0
       \quad(0\ne\lambda\in\Lambda_G(E)).
\]
The quantitative descent criterion of
\cite[Section 5.4]{ArithmeticReductions} gives finite
generation and fixed-field Northcott. If \(B\ge0\),
\(j\in\Z_{\ge0}\), and \(m^{2j}\varepsilon_E>4B\), then
\[
 \#\{\lambda\in\Lambda_G(E):q_G(\lambda)\le B\}\le s^j.
\]
The reason is that two classes in the ball congruent
modulo \(m^j\Lambda_G(E)\) would differ by \(m^jz\) with
\(q_G(z)<\varepsilon_E\). The usual height descent then
generates the group from a finite ball and representatives
modulo \(m\).

Scalar detection at one prime supplies a finite quotient
with \(m=p\), through the local integral bound.
A genuine positive height gap is the additional input.
Real positive definiteness alone does not supply this gap
on a subgroup with unbounded denominators.

\section{Minimal class fields and bounded-degree Northcott}
\label{sec:class-fields}

\subsection{Trace and the class field}\label{sec:class-trace}

View the residual groups inside the geometric quotient
with its discrete \(\Gal(\Fbar/F_0)\)-action. For a class
\(\xi\), define
\[
 F_0(\xi)=
 \Fbar^{\operatorname{Stab}_{\Gal(\Fbar/F_0)}(\xi)}.
\]
If \(\xi\in\Lambda_G(E)\), its minimal class field
\(K=F_0(\xi)\) is contained in \(E\), and trace gives
\begin{equation}\label{eq:class-trace}
 [E:K]\xi\in\Lambda_G(K).
\end{equation}
Indeed the pushforward of an integral representative has
geometric class equal to the sum over the \(K\)-embeddings
of \(E\), and every summand is \(\xi\). Thus the multiplier
is at most \(D\) when \([E:F_0]\le D\).
This concerns the field of the class; a particular cycle
representative need not descend to \(K\).

\begin{conjecture}[Residual class-field ramification]
\label[conjecture]{conj:ramification}
For every \(B\ge0\) and integer \(D\ge1\), there is a
finite set \(S(B,D)\) of finite places of \(F_0\) such that
the fields \(F_0(\xi)/F_0\), for
\[
 \xi\in\Lambda_G(E),\qquad [E:F_0]\le D,\qquad q_G(\xi)\le B,
\]
are unramified outside \(S(B,D)\).
\end{conjecture}

Under fixed-field Northcott for every finite extension,
\Cref{conj:ramification} is equivalent to bounded-degree
Northcott. Hermite--Minkowski leaves finitely many possible
class fields \(K\) of degree at most \(D\).
Equation \eqref{eq:class-trace} then reduces the classes to
the finitely many fixed-field balls
\[
 m\xi\in\Lambda_G(K),\qquad
 1\le m\le D,\qquad q_G(m\xi)\le D^2B.
\]
The converse follows by taking the union of ramification
sets of the finitely many classes in a ball.
The full proof is in
\cite[Section 6.2]{ArithmeticReductions}; the
finite-ramification form of Hermite--Minkowski is
\cite[Theorem 10.2.2]{Kedlaya}.

\subsection{The extra information needed as the field varies}
\label{sec:varying-field-obstacle}

A fixed set \(S_E\) for \(H^1(E,V_p)\) controls ramification
over \(E\). It does not bound the ramification of the
class field \(F_0(\xi)/F_0\). The distinction persists even
when fixed-field estimates are uniform in the degree.
For example, choose infinitely many distinct quadratic
extensions \(K_i/F_0\), with characters \(\chi_i\), and set
\[
 \mathcal G=\bigoplus_i\Q e_i,\qquad
 \sigma e_i=\chi_i(\sigma)e_i,\qquad
 \Lambda(E)=\bigoplus_{K_i\subseteq E}\Z e_i,\qquad
 q\left(\sum a_ie_i\right)=\sum a_i^2.
\]
These integral images are compatible with trace.
Their ranks, height gaps and quotients modulo any integer
\(m\ge2\) have degree-dependent uniform bounds, and all
fixed-field balls are finite. Yet the infinitely many
\(e_i\in\Lambda(K_i)\) have degree two and height one.
This example isolates the field-control hypothesis in
the ramification criterion.

A quantitative strengthening asks for constants
\(A_D,B_D\ge0\), depending only on the fixed residual data
and \(D\), such that
\begin{equation}\label{eq:discriminant-target}
 \log N_{F_0/\Q}\mathfrak d_{F_0(\xi)/F_0}
    \le A_Dq_G(\xi)+B_D
 \qquad([E:F_0]\le D,\ \xi\in\Lambda_G(E)).
\end{equation}
This is a research question stronger than
\Cref{conj:ramification}. A bound of the displayed form
would imply the required finite ramification set at every
bounded height.

\section{A proof programme and its arithmetic obstacles}
\label{sec:programme}

The reduction theorem provides two complementary routes.
One seeks the three central assertions directly. The other
constructs arithmetic detectors for the residual quotient
and then proves integral and field-of-definition control.

\subsection{The geometric and central steps}\label{sec:geometric-steps}

The first target is \Cref{conj:offcentral}. The tensor Chow
decomposition gives an exact test: a nonzero off-central
core would produce a misplaced Lefschetz string.
A proof must therefore eliminate the core in the rational
Chow group, rather than merely its cohomological realization.

The next target is the algebraic kernel in
\Cref{conj:AJ}. The projector construction already identifies
it as \(R_SA_{S,E}\), and adjunction places it in the full
central height radical. Establishing its vanishing is a
separate arithmetic cycle statement. Once it vanishes,
the visible Néron--Tate factor has all the required
fixed-field and bounded-degree finiteness properties.

For the residual factor, real positivity is an initial
target. It supplies Cauchy--Schwarz and, together with the
algebraic kernel statement, central nondegeneracy.
The primitive sum-of-squares identity transfers these
consequences uniformly to every degree of the product.

\subsection{A detector-to-Northcott route}\label{sec:arithmetic-steps}

For each fixed finite \(E/F_0\), the sequence of inputs is:
\begin{enumerate}
\item Prove scalar-extended injectivity of one residual
regulator. The finite-dimensional cohomological target then
gives finite rational rank.
\item Prove one common denominator bound for the integral
image, for example through the all-prime saturation
condition \eqref{eq:common-lattice-target}. Alternatively,
combine a finite Kummer quotient with a positive height gap.
\item Establish real residual positivity. Together with
the common lattice, this gives fixed-field Northcott.
\item Control the ramification of minimal class fields in
bounded degree and height. The trace criterion then gives
bounded-degree Northcott.
\end{enumerate}
The alternatives in the second step have different
arithmetic content. All-prime saturation compares integral
cohomology across coefficient primes. The height-gap
route instead requires a uniform positive lower bound on
nonzero integral classes over the fixed field.

Under real positivity, scalar detection and the
almost-everywhere lattice comparison for every finite
\(E/F_0\), the first three steps establish fixed-field
Northcott. Adding \Cref{conj:ramification} proves
\Cref{conj:residual-northcott}. Alternatively, a finite
Kummer quotient and a genuine gap over every fixed field,
together with real positivity and
\Cref{conj:ramification}, give the same conclusion without
a preliminary finite-rank assumption. Finally,
\Cref{conj:offcentral,conj:AJ} and
\Cref{thm:program-equivalence} give the product programme.

The local-to-global construction of a height system and
these arithmetic properties meet at \Cref{def:heights}.
A construction of local Green forms must supply the
global pairing, adjunction and curve normalization before
the arithmetic reductions apply. The present programme
starts from that explicitly stated height system; its
subsequent tasks concern vanishing, positivity, integral
discreteness and fields of cycle classes.


% Source: papers/arithmetic-sections/conjecture-references.tex
\begin{thebibliography}{99}
\addcontentsline{toc}{section}{References}

\bibitem{ArithmeticReductions}
\emph{Arithmetic Lefschetz and integral heights:
proofs and reductions on products of curves}, companion research
manuscript, 17 November 2026.

\bibitem{Kahn}
B. Kahn, \emph{On the injectivity and non-injectivity of the
\(\ell\)-adic cycle class maps}, manuscript dated 10 March 2024.
\href{https://webusers.imj-prg.fr/~bruno.kahn/preprints/ladic-cycle-class-rev3.pdf}
{Author-hosted PDF}.

\bibitem{Kedlaya}
K. S. Kedlaya, \emph{Weil cohomology in practice},
Chapter 10, Comparing Galois representations: the Faltings--Serre
method, lecture of 30 October 2019.
\href{https://kskedlaya.org/weil-cohom/chapter-10.html}
{Author-hosted notes}.

\bibitem{Raskind}
W. Raskind, \emph{Higher \(\ell\)-adic Abel--Jacobi mappings and
filtrations on Chow groups},
\href{https://arxiv.org/abs/alg-geom/9412001}{arXiv:alg-geom/9412001}.

\bibitem{ScaviaSuzuki}
F. Scavia and F. Suzuki,
\emph{Non-injectivity of the cycle class map in continuous
\(\ell\)-adic cohomology},
\href{https://arxiv.org/abs/2208.06713}{arXiv:2208.06713}.

\bibitem{ZhangStandard}
S.-W. Zhang, \emph{Standard conjectures and height pairings},
\href{https://arxiv.org/abs/2009.07089v3}{arXiv:2009.07089v3}.
The numbering cited here refers to the version of 29 April 2022.

\end{thebibliography}



\end{document}

